Note~33 makes the funding distinction measurable for 2025. Its parent-only
income statement reports net income of US\$57{,}048m. The parent-only cash-flow
statement subtracts US\$71{,}163m of subsidiary and affiliate net income,
reports a US\$14{,}115m parent company net loss in that reconciliation, then
adds US\$50{,}000m of cash dividends from subsidiaries and US\$8{,}583m of
other operating adjustments to reach US\$44{,}468m of operating cash. Neither
the income figure nor the dividend alone is the period's ending cash.

At 31~December 2025, the parent reports US\$39m of cash and due from banks and
US\$9{,}751m of deposits with banking subsidiaries, together US\$9{,}790m.
The cash-flow footnote says this measure includes US\$8.0bn of intercompany
time deposits. The same parent-only balance sheet reports US\$669{,}449m of
equity-method investments in bank and bank-holding subsidiaries; that asset is
not parent cash. Operating cash of US\$44{,}468m and investing cash of
US\$34m, less US\$44{,}512m used in financing, leave closing cash and
deposits US\$10m below their opening balance.
These dated amounts distinguish accounting income, existing cash and
conditional subsidiary transfers. A future debt-service calculation still
needs payment dates, permissible transfers, deposit terms and funding actions
\citep[pp.~19, 310--311]{jpmorgan2025tenk}.

\newpage

